\documentclass[11pt]{article}
\usepackage[margin=1in]{geometry}
\usepackage[T1]{fontenc}
\usepackage[utf8]{inputenc}
\usepackage{lmodern}
\usepackage{hyperref}
\usepackage{booktabs}
\usepackage{enumitem}

\title{A Ledger-Aware Benchmark for Market Infrastructure under Frequent Batch Auctions}
\author{Author Name \\ Institution \\ \texttt{email@example.com}}
\date{\today}

\begin{document}
\maketitle

\begin{abstract}
We present a benchmark-oriented extension of a ledger-first market infrastructure prototype for studying market-design and execution behavior under settlement and risk constraints. The environment combines seedable agent-based order flow, configurable immediate and frequent-batch-auction matching regimes, double-entry style account state transitions, and deterministic safety checks for conservation and non-negativity. We expose reproducible benchmark tasks that compare latency, throughput, spread, price impact, queue-priority advantage, and latency-arbitrage profit across market designs.
\end{abstract}

\section{Introduction}
This paper is a parallel benchmark track built on top of an existing market-infrastructure systems prototype. The objective is not to replace the original systems paper, but to provide a reusable evaluation environment for market-design and execution-policy studies under ledger-aware settlement constraints.

\section{Research Questions}
We focus on three questions:
\begin{enumerate}[leftmargin=*]
  \item How do immediate and frequent-batch-auction regimes differ in latency and fill behavior?
  \item Can fairness-adjacent signals such as queue-priority advantage and latency-arbitrage profit be measured reproducibly?
  \item Do settlement invariants hold under heterogeneous agent populations and stress scenarios?
\end{enumerate}

\section{Environment}
The benchmark environment consists of seedable agent models, configurable matching modes, deterministic settlement application, and generated benchmark artifacts. The current implementation covers immediate execution, fixed batch windows, and a stress scenario.

\section{Metrics}
We report:
\begin{itemize}[leftmargin=*]
  \item orders/s and fills/s
  \item p50/p95/p99 latency
  \item average spread and price impact
  \item queue-priority advantage, execution dispersion, and arbitrage profit
  \item settlement safety violations
\end{itemize}

\section{Current Results}
\begin{table}[h]
\centering
\begin{tabular}{lrrrr}
\toprule
Scenario & Orders/s & Fills/s & p50 (ms) & p99 (ms) \\
\midrule
Immediate-Surrogate & 1360.0 & 829.2 & 10 & 10 \\
FBA-100ms & 1360.0 & 802.5 & 50 & 310 \\
FBA-250ms & 1358.4 & 628.8 & 80 & 490 \\
FBA-500ms & 1358.7 & 644.7 & 190 & 910 \\
FBA-250ms-Stress & 1775.2 & 987.2 & 100 & 590 \\
\bottomrule
\end{tabular}
\caption{Current simulator benchmark snapshot. All scenarios report zero conservation breaches and zero negative-balance violations.}
\end{table}

\section{Limitations}
The benchmark currently uses fixed policies and proxy fairness metrics. It does not yet provide multi-seed aggregate reporting, adaptive execution baselines, or richer market-impact simulation.

\section{Conclusion}
This benchmark track establishes a reproducible, ledger-aware evaluation environment for market-design experiments without overwriting the original systems-paper line.

\bibliographystyle{plain}
\bibliography{references}

\end{document}
